\documentclass[twoside]{article}
\usepackage[utf8]{inputenc}

% The package supports other formats as well, but this template
% is primarily checked and updated for examinations at TUM BGDM
%\usepackage[uebung, answers]{tumbgdm}

\usepackage[uebung]{tumbgdm}
\usepackage{tabularx}
\usepackage{amsmath}

\Professor{Prof. Dr. Martin Werner}
\Universitaet{Technical University of Munich}
\Institut{School of Engineering and Design}
\Lehrstuhl{Professorship Big Geospatial Data Management}
\Semester[WiTe 2024/25]{Winter Term 2024/25}
\Vorlesung{Principles of Programming}
\InstitutsWebsite{https://www.bgd.lrg.tum.de/}
\VorlesungsWebsite{https://www.moodle.tum.de/course/view.php?id=103704}

\nummerUebungsblatt{11}
\nummerErsteFrage{1}

\LearningOutcome{ Exam Preparation }

\begin{document}
\maketitle

\vspace{10pt}

In this tutorial, we practice a mock exam.

\vspace{10pt}

\begin{task}{State your Name}{}{}

Name: \\
 \\
 \\
Surname: \\
 \\
 \\

\end{task}

\begin{task}{Yes-No-Example}{}{}
Answer the following questions with yes or no. \\

\begin{yesnolist}
  \yesnoquestion{In Python, you need to define a function to execute a script.}{X}{}{}
  \yesnoquestion{The Numpy and Matplotlib library are by default part of the Python installation.}{X}{}{}
  \yesnoquestion{To install a Python package, you always have to create a virtual environment.}{X}{}{} 
\end{yesnolist}
\end{task}

\clearpage

\begin{task}{Multiple Joice Example}{}{}
For the given question, one of the answer options is correct. Mark this (and only this) answer. If more than one answer is marked, the answer is not considered. A wrong answer is \textbf{not} counted as -1.\vspace{1cm}

  \begin{multiplechoice}{In programming, what is meant by a recursive function?}
    \mcanswer{A function that is called multiple times.}{}
    \mcanswer{A function that calls itself.}{}
    \mcanswer{A function inside another function.}{}
    \mcanswer{A function that takes another function as input.}{}
  \end{multiplechoice}

  \begin{multiplechoice}{Which of the following is \textbf{not} a method of a list object in Python?}
    \mcanswer{remove}{X}
    \mcanswer{delete}{}
    \mcanswer{reverse}{}
    \mcanswer{pop}{}
  \end{multiplechoice}
\end{task}


\begin{task}{Distance Between Two Points}{}{}
Complete the given Python function \code{dist()} that takes two points \code{p1 = [x1, y1]} and \code{p2 = [x2, y2]} as input and shall return their distance.
\begin{lstlisting}
    def dist(p1, p2):
        # ...
        return distance
\end{lstlisting}
\end{task}

\clearpage

\begin{task}{List Comprehension}{}{}
  \begin{enumerate}
  \item{Consider the following program using list expressions.}
  \begin{lstlisting}
  vocabulary = ['dear', 'done', 'easy', 'hello', 'hope', 'is', 'it', 'say', 'student', 'to', 'we', 'well']
  sentence = " ".join([vocabulary[x] for x in [3,0,8,10,4,6,5,2]])
  print(sentence)
  \end{lstlisting}
  Give the exact output of the program 

\vspace*{6cm}


  \item{Consider the following Python function consisting of seven lines of code.}
  \begin{lstlisting}
  def myFunction(n: list):
      result = []
      for k in n:
          if k > 0:
              tmp = k**2
              result.append(tmp)
      return result
  \end{lstlisting}
  Rewrite the function with a list comprehension such that it only needs two lines of code.
  \end{enumerate}

\end{task}

\end{document}